% !TeX root = ../../Book3.tex
\chapter{Ahlfors–Schwarz 引理与 Picard 定理}
\label{b3:ch:14}
\BookChapterNumber{b3:ch:14}
\begin{chapterabstract}
本章完整证明 Ahlfors--Schwarz 曲率比较，并把它与覆盖双曲几何相连。球面导数、Marty 判别和重标度给出省略值 Montel 定理，进而完整推出 Picard 小、大定理。选读部分明确区分统一负曲率上界与单纯点态负曲率，并解释紧曲面上的 Brody 思想。
\end{chapterabstract}

\begin{learninggoals}
\begin{itemize}
\item\label{b3:item:14-goal-1}计算共形曲率与拉回密度，正确处理临界点。
\item\label{b3:item:14-goal-2}用内部最大值证明负曲率密度比较及其完备源形式。
\item\label{b3:item:14-goal-3}证明 Marty、重标度与省略值正规性，辨认依赖的覆盖输入。
\item\label{b3:item:14-goal-4}完整推出 Picard 小、大定理并核对例外值及无限原像量词。
\end{itemize}
\end{learninggoals}

Schwarz 引理曾把一个内点归一化变为刚性。曲率方法则用整个密度的二阶不等式代替自同构归一化。省略值把目标变成双曲穿孔球面，球面紧性又把局部异常导数变成整个平面极限，最后导出 Picard 理论。

\ProseParagraphBreak
主线依赖第13章明确列出的覆盖模型输入、以及此前的 Hurwitz、Liouville 和 Casorati--Weierstrass。Ahlfors--Schwarz、Marty、重标度和两项 Picard 证明均在本章完成；一般统一化的深层证明状态不会因此被误记为已完成。

% !TeX root = ../../Book3.tex
\section{共形度量}
\label{b3:sec:1401}

第十三章的 Poincaré 密度同时具备三种性质：它给出路径长度，在全纯坐标之间按导数的模变换，并且具有常负曲率。本章保留前两种结构，把常曲率条件放宽为一个曲率不等式。为此，先说明一般密度怎样给出距离、拉回在临界点发生什么，以及完备性究竟要求控制哪些路径。

\subsection{共形度量}
\label{b3:subsec:140101}

设 \(U\subseteq\mathbb C\) 是区域。正函数 \(\rho:U\to(0,\infty)\) 定义长度元
\[
 \mathrm ds_\rho=\rho(z)|\mathrm dz|.
\]
它在每一点把所有方向的欧氏长度乘以同一个因子，因而保持夹角，称为共形长度元。其面积元是 \(\rho(z)^2\,\mathrm dA(z)\)，因为两个相互垂直方向的长度都乘以 \(\rho(z)\)。讨论长度与拓扑只需 \(\rho\) 连续；计算曲率时则要求 \(\rho\) 为 \(C^2\) 函数。

\ProseParagraphBreak
若 \(w=F(z)\) 是全纯坐标变换，\(F'\ne0\)，同一长度元在两种坐标中的密度须满足
\[
 \rho_z(z)=\rho_w(F(z))\cdot|F'(z)|.
\]
因此，Riemann 曲面上的一个共形度量，是每张全纯坐标图中的正密度，且这些密度在重叠处满足此式。沿分段 \(C^1\) 路径 \(\gamma\) 定义
\[
 \ell_\rho(\gamma)=\int\rho(\gamma(t))\cdot|\gamma'(t)|\,\mathrm dt.
\]
曲面上使用有限个坐标段分别积分；变换公式和链式法则保证各段的数值不依赖坐标，进一步细分也不改变总长度。以全部连接路径的长度下确界定义 \(\mathrm d_\rho(p,q)\)。对于正密度，这确实给出曲面的距离，但证明正定性时还要考虑路径离开坐标盘的可能。

\Needspace{6\baselineskip}
\begin{proposition}[正密度的路径距离]\label{b3:prop:14-positive-density-distance}
设 \(X\) 是连通 Riemann 曲面，\(\rho\) 是 \(X\) 上正连续共形密度。由分段 \(C^1\) 路径长度的下确界定义的 \(\mathrm d_\rho\) 是有限值距离，其拓扑等于 \(X\) 原有的拓扑。具体地，固定 \(p\in X\)，在一个坐标中令 \(p\) 对应 \(a\)，并取紧包含于该坐标图的闭圆盘 \(\overline{D(a,r)}\)。若该盘上的密度满足 \(0<m\le\rho\le M\)，则对盘内坐标为 \(z\) 的点 \(q\)，有
\[
 m|z-a|\le\mathrm d_\rho(p,q)\le M|z-a|.
\]
对不在开坐标盘 \(D(a,r)\) 内的点 \(q\)，还有 \(\mathrm d_\rho(p,q)\ge mr\)。
\end{proposition}
\begin{proof}
坐标内的直线段给出上界。对于任意连接 \(p,q\) 的路径，若它一直留在该坐标盘中，则密度下界和欧氏曲线长度下界给出 \(\ell_\rho(\gamma)\ge m|z-a|\)。若路径离开该盘，取第一次到达圆周的初始路径段。该段的欧氏长度至少为 \(r\)，故整个路径长度至少为 \(mr\ge m|z-a|\)。这证明盘内点的下界，也同时证明盘外点的下界。

连通曲面上的任意两点可用有限串相交坐标盘及盘内折线路径连接；每段路径的像紧，连续密度在其上有界，所以距离有限。上述局部下界保证不同点的距离为正。常路径、反向路径和路径拼接分别给出零对角线、对称性与三角不等式；三角不等式的下确界步骤与命题~\ref{b3:prop:13-disk-path-metric} 相同，只需分别选取长度在距离的 \(\varepsilon/2\) 以内的两条路径。

欧氏坐标收敛由局部上界推出距离收敛。反之，若 \(\mathrm d_\rho(p,q_n)\to0\)，盘外下界保证 \(q_n\) 最终留在上述坐标盘中，盘内下界再保证其坐标趋于 \(a\)。两种拓扑因此相同。
\end{proof}

拉回全纯映射时还会遇到非负密度。这里称连续非负的共形长度元为共形伪度量，允许它在某些点退化；若其零点离散，则在零点之外仍可按正度量计算。这个用语指长度元允许为零，并不预先保证由它定义的路径下确界能分离各点。后面的曲率比较只使用密度的连续性及其正值处的二阶微分，不在零点定义 \(\log\rho\)。

\subsection{Gaussian 曲率}
\label{b3:subsec:140102}

对正 \(C^2\) 共形密度，本章沿用曲率表达
\[
 K_\rho\coloneqq-\frac{\Delta\log\rho}{\rho^2},
 \quad \Delta=\partial_x^2+\partial_y^2.
\]
它与二维几何中 Gaussian 曲率的认同，按第十三章所说明的范围作为几何输入\footnote{Ahlfors《Conformal Invariants》\cite[1-5 Ultrahyperbolic Metrics, p.~12]{Ahlfors2010Conformal}使用这一表达讨论共形曲率。本书需要的坐标不变性已在命题~\ref{b3:prop:13-curvature-coordinate} 中完整计算；曲率与一般内在几何定义的认同不另假定为本书已证明的曲率张量理论。}。下面恢复本章使用它时的条件与符号。分母记录面积尺度，分子记录密度对数的二阶变化；常密度的曲率因而为零。

\Needspace{6\baselineskip}
\begin{proposition}\label{b3:prop:14-curvature-formula}
设 \(X\) 是 Riemann 曲面，\(\rho\) 是其正 \(C^2\) 共形密度。在任意局部全纯坐标中定义
\[
 K_\rho=-\rho^{-2}\Delta\log\rho,
 \quad \Delta=\partial_x^2+\partial_y^2.
\]
则这一表达给出 \(X\) 上与坐标无关的函数。对常数 \(\kappa>0\)，条件 \(K_\rho\le-\kappa\) 等价于各坐标中均有
\[
 \Delta\log\rho\ge\kappa\rho^2.
\]
\end{proposition}
\begin{proof}
在两张重叠坐标之间，设 \(w=F(z)\)，则 \(F'\ne0\)，密度关系为 \(\rho_z=(\rho_w\circ F)\cdot|F'|\)。命题~\ref{b3:prop:13-curvature-coordinate} 已在正 \(C^2\) 密度及非零全纯坐标导数的条件下证明
\[
 -\frac{\Delta_z\log\rho_z}{\rho_z^2}
 =\left(-\frac{\Delta_w\log\rho_w}{\rho_w^2}\right)\circ F.
\]
该命题使用的两项恒等式是 \(\Delta_z(v\circ F)=|F'|^2(\Delta_wv)\circ F\) 与 \(\Delta\log|F'|=0\)，此处的坐标与密度满足全部条件，故公式直接适用。最后，将 \(-\rho^{-2}\Delta\log\rho\le-\kappa\) 乘以正数 \(\rho^2\) 并移项，得到所述等价不等式。
\end{proof}

对径向函数 \(v=v(r)\)，在 \(r>0\) 处有 \(\Delta v=v''+v'/r\)。将 \(v=\log(1\pm r^2)\) 代入，并在原点按连续性延拓，得到
\[
 \Delta\log\frac{2}{1-|z|^2}
 =\frac{4}{(1-|z|^2)^2},
 \quad
 \Delta\log\frac{2}{1+|z|^2}
 =-\frac{4}{(1+|z|^2)^2}.
\]
连同 \(\Delta\log1=0\)，即得\cref{b3:tab:14-curvature-models} 中的三种曲率。球面密度在平面坐标中趋于零，但这只是该坐标接近无穷远点的表现：换成 \(\zeta=1/z\) 后，密度重新成为 \(2/(1+|\zeta|^2)\)，在 \(\zeta=0\) 处光滑且为正。

\begin{table}[htbp]
\centering\small\renewcommand{\arraystretch}{1.18}\setlength{\tabcolsep}{4pt}
\caption[三种标准共形曲率]{三种标准共形曲率。球面行采用球面去掉一点后的平面坐标；补上无穷远点才得到整张球面。}
\label{b3:tab:14-curvature-models}
\begin{tabularx}{\linewidth}{@{}>{\raggedright\arraybackslash}p{2cm}>{\centering\arraybackslash}p{3.3cm}>{\centering\arraybackslash}p{1.3cm}>{\raggedright\arraybackslash}X@{}}
\toprule
模型 & 密度 \(\rho(z)\) & 曲率 \(K\) & 坐标定义域 \\
\midrule
Euclid 平面 & \(1\) & \(0\) & \(\mathbb C\) \\
Poincaré 圆盘 & \(2/(1-|z|^2)\) & \(-1\) & \(|z|<1\) \\
单位球面 & \(2/(1+|z|^2)\) & \(1\) & \(\mathbb C\)；另一坐标补无穷远点 \\
\bottomrule
\end{tabularx}
\end{table}

若把密度乘以常数 \(c>0\)，则
\[
 \ell_{c\rho}=c\ell_\rho,
 \quad \mathrm d_{c\rho}=c\,\mathrm d_\rho,
 \quad K_{c\rho}=c^{-2}K_\rho.
\]
前两式来自长度积分，第三式来自 \(\log(c\rho)=\log c+\log\rho\)。面积同时乘以 \(c^2\)，而 Cauchy 列与完备性不变。例如，密度 \(1/(1-|z|^2)\) 的曲率是 \(-4\)；将它乘以 \(2\)，才得到本卷使用的曲率 \(-1\) 的圆盘度量。曲率比较中的常数因此必须与长度归一化一起保留。

\ProseParagraphBreak
局部曲率还不能辨认全局拓扑。去心平面上的密度 \(\rho(z)=1/|z|\) 满足 \(\Delta\log\rho=0\)，因为 \(\log|z|\) 在 \(z\ne0\) 处调和。在任一有对数分支的邻域中，\(\mathrm dw=\mathrm dz/z\)，故这个长度元就是 \(w=\log z\) 坐标中的 Euclid 长度元；绕原点一周后 \(w\) 增加 \(2\pi\mathrm i\)，于是全局模型是平直圆柱。平面与圆柱的曲率都为零，闭路和覆盖资料却不同。

\subsection{拉回度量}
\label{b3:subsec:140103}

设 \(f:X\to Y\) 是 Riemann 曲面之间的全纯映射，目标有正 \(C^2\) 共形密度 \(\rho\)。在源坐标 \(z\) 与目标坐标 \(w\) 中，令 \(F=w\circ f\circ z^{-1}\) 是 \(f\) 的局部表达，定义拉回密度
\[
 h_z(z)=\rho_w(F(z))\cdot|F'(z)|.
\]
换目标坐标时，\(\rho_w\) 与 \(F'\) 中的坐标变换导数相消；换源坐标时，\(h_z\) 正好按共形密度的规律变换。因此它定义了 \(X\) 上的全局非负长度元，记为 \(f^*\rho\)。这个记号指拉回后的密度系数，而每个实际计算仍在局部坐标中进行。

\Needspace{6\baselineskip}
\begin{proposition}\label{b3:prop:14-pullback-curvature}
设 \(X,Y\) 是连通 Riemann 曲面，\(f:X\to Y\) 非常值全纯，\(\rho\) 是 \(Y\) 上正 \(C^2\) 共形密度。令 \(h=f^*\rho\)，即局部坐标中的 \(h=(\rho\circ f)\cdot|f'|\)。则 \(h\) 连续，其零点恰为 \(f\) 的临界点，且这些零点离散。在非临界点处，\(h\) 为正 \(C^2\) 密度，并满足
\[
 K_h=K_\rho\circ f.
\]
对任意分段 \(C^1\) 源路径 \(\gamma\)，包括经过临界点的路径，均有
\[
 \ell_h(\gamma)=\ell_\rho(f\circ\gamma).
\]
\end{proposition}
\begin{proof}
局部表达 \(F\) 非常值，故 \(F'\) 不恒为零，其零点离散。目标密度严格为正，所以 \(h\) 的零点正是 \(F'\) 的零点。坐标变换导数不为零，这个零点集合不依赖坐标。连续性由拉回公式得到。

在 \(F'\ne0\) 的点附近，\(F\) 是全纯坐标变换，命题~\ref{b3:prop:14-curvature-formula} 的变换式直接给出 \(K_h=K_\rho\circ f\)。在任意坐标路径段上，链式法则给出
\[
 \rho(F(\gamma(t)))\cdot|(F\circ\gamma)'(t)|
 =\rho(F(\gamma(t)))\cdot|F'(\gamma(t))|\cdot|\gamma'(t)|.
\]
这个等式在 \(F'=0\) 时仍成立；积分并将有限个坐标段相加，便得长度等式。
\end{proof}

临界点附近的退化有明确形式。若 \(F'\) 在 \(a\) 处有 \(m\) 阶零点，写成 \(F'(z)=(z-a)^m g(z)\)，其中 \(g\) 全纯且在邻域内无零点，则
\[
 h(z)=|z-a|^m b(z),
 \quad b(z)=\rho(F(z))\cdot|g(z)|>0,
\]
且 \(b\) 为 \(C^2\) 函数。因子 \(|z-a|^m\) 说明拉回密度在临界点可能不再是正 \(C^2\) 密度；曲率公式不能在那里除以 \(h^2\)。若 \(f\) 为常值，整个拉回密度都为零。下一节将允许这些退化，并只在密度为正的比较点计算曲率。

\subsection{完备性}
\label{b3:subsec:140104}

正连续密度 \(\rho\) 称为完备的，是指其路径距离 \(\mathrm d_\rho\) 的每个 Cauchy 列都在曲面内收敛。正性保证局部距离与原拓扑一致，完备性则限制路径能否用有限长度离开曲面的全部紧部分。为了精确表达后一点，称连续路径 \(\gamma:[0,1)\to X\) 逃离紧集，是指对任意紧集 \(K\subseteq X\)，都存在 \(t_K<1\)，使 \(\gamma(t)\notin K\) 对所有 \(t>t_K\) 成立。

\Needspace{6\baselineskip}
\begin{proposition}[完备性的逃逸路径判据]\label{b3:prop:14-completeness-escape}
设 \(X\) 是连通 Riemann 曲面，\(\rho\) 是其正连续共形密度，\(\mathrm d_\rho\) 是相应路径距离。则 \((X,\mathrm d_\rho)\) 完备，当且仅当每条逃离紧集的连续、局部分段 \(C^1\) 路径 \(\gamma:[0,1)\to X\) 都具有无限长度。这里总长度定义为
\[
 \ell_\rho(\gamma)=\sup_{0<t<1}\ell_\rho(\gamma|_{[0,t]}).
\]
\end{proposition}
\begin{proof}
先设距离完备。若有逃离紧集的路径长度有限，则其尾部长度趋于零。取 \(t_n\uparrow1\)，对 \(m>n\) 有
\[
 \mathrm d_\rho(\gamma(t_n),\gamma(t_m))
 \le\ell_\rho(\gamma|_{[t_n,t_m]})\longrightarrow0
\]
且上界对所有 \(m>n\) 一致。因此 \(\gamma(t_n)\) 是 Cauchy 列，收敛到某个 \(p\in X\)。由距离拓扑与曲面拓扑相同，取 \(p\) 的一个紧坐标邻域后，该列最终留在此紧集中，与路径逃离每个紧集矛盾。

反之，若距离不完备，取一个没有极限的 Cauchy 列。可抽取子列 \((p_n)\)，使
\[
 \mathrm d_\rho(p_n,p_{n+1})<2^{-n-2}.
\]
按路径下确界的定义，选取连接 \(p_n,p_{n+1}\) 的分段 \(C^1\) 路径 \(\gamma_n\)，使 \(\ell_\rho(\gamma_n)<2^{-n-1}\)。把这些路径依次线性改参数到区间 \([1-2^{1-n},1-2^{-n}]\)，拼成 \(\gamma:[0,1)\to X\)。它连续且局部分段 \(C^1\)，总长度不超过 \(\sum_{n\ge1}2^{-n-1}<\infty\)。

还须证明这条路径确实逃离紧集。若某个紧集 \(K\) 被任意靠后的路径段访问，就能选取 \(n_j\to\infty\) 和 \(q_j\in K\cap\gamma_{n_j}\)。由紧性及距离拓扑与原拓扑相同，可再取子列使 \(q_j\to q\in K\)。而
\[
 \mathrm d_\rho(p_{n_j},q_j)
 \le\ell_\rho(\gamma_{n_j})<2^{-n_j-1}\longrightarrow0,
\]
故 \(p_{n_j}\to q\)。原来的 Cauchy 列一旦有收敛子列，就必须整体收敛到同一点，矛盾。因此 \(\gamma\) 逃离每个紧集，构成有限长度的逃逸路径。这证明逆命题。
\end{proof}

这个判据也给出实用的比较办法。若 \(\sigma\) 是同一曲面上的完备正连续密度，而 \(\rho\ge c\sigma\)、\(c>0\)，则每条逃逸路径都有
\(\ell_\rho\ge c\ell_\sigma=\infty\)，所以 \(\rho\) 完备。例如，平面密度 \(e^{|z|^2}\ge1\)，由 Euclid 平面的完备性即可知道它完备；无需先寻找该度量的最短路径。

\ProseParagraphBreak
单位盘的常密度 \(1\) 不完备，因为径向路径可在有限长度内趋近单位圆周。圆盘 Poincaré 密度的完备性已在第十三章证明；若将同一密度限制到 \(D(0,a)\)、\(0<a<1\)，曲率仍恒为 \(-1\)，但趋近实轴端点 \(a\) 的径向长度只有 \(2\operatorname{arctanh}a\)，所以这个受限度量不完备。负曲率条件由局部微分表达控制，端部是否在有限距离处还须另行判断。

\ProseParagraphBreak
表中的球面密度也说明不能忘记曲面本身的范围。在 \(\mathbb C\) 上，\(\rho(z)=2/(1+|z|^2)\) 沿正实轴到无穷远的总长度为 \(\int_0^\infty2/(1+t^2)\,\mathrm dt=\pi\)，故去掉无穷远点的度量不完备；补上该点后，密度在另一坐标中仍正且光滑，所得球面是紧度量空间，因而完备。

\ProseParagraphBreak
应用逃逸判据时，一条有限长度的逃逸路径就足以否定完备性；证明完备则必须控制所有逃逸路径。只算几条径向线的长度不能承担后一种量词。下一节的 Ahlfors--Schwarz 比较直接使用局部曲率和小圆盘最大值，不把完备性列为待比较密度的假设。

% !TeX root = ../../Book3.tex
\section{Ahlfors–Schwarz 引理}
\label{b3:sec:1402}

Ahlfors--Schwarz 比较使用一个在边界发散的圆盘密度，将正的最大值强迫到内部，再由 Laplace 不等式排除它。这个证明只需二阶微分与最大值原理。

\subsection{曲率比较}
\label{b3:subsec:140201}

\term[ahlfors-schwarz-yinli]{Ahlfors--Schwarz 引理}[Ahlfors--Schwarz lemma]用统一负曲率界比较密度。

\begin{theorem}[Ahlfors--Schwarz]\label{b3:thm:14-ahlfors-schwarz}
若 \(\rho>0\) 为 \(D(0,R)\) 上 \(C^2\) 密度，\(K_\rho\le-\kappa<0\)，则
\[
 \rho(z)\le\frac{2R}{\sqrt\kappa(R^2-|z|^2)}.
\]
\end{theorem}
\begin{proof}
固定 \(r<R\)，令 \(\eta_r=2r/(\sqrt\kappa(r^2-|z|^2))\)。它满足 \(\Delta\log\eta_r=\kappa\eta_r^2\)。\(\rho\) 在闭半径 \(r\) 盘有界，\(\eta_r\) 在边界发散，所以 \(\rho/\eta_r\rightarrow0\)。若比值曾大于1，\(v=\log(\rho/\eta_r)\) 有内部正最大值。在该点
\[
 \Delta v\ge\kappa(\rho^2-\eta_r^2)>0,
\]
与最大值处 \(\Delta v\le0\) 矛盾。故 \(\rho\le\eta_r\)，令 \(r\uparrow R\) 即得结论。
\end{proof}
固定较小圆盘再令半径增大，是为了只使用内部有界性，不假定原密度有边界延拓。等号密度本身在最外边界无界，并且准确达到比较常数。

\ProseParagraphBreak
比较证明先用半径 \(r<R\) 的模型，是为创造一个可达到的内部最大值。模型密度在 \(|z|=r\) 发散，而给定密度在闭 \(r\) 盘附近有界，所以比值趋于0，正的最大比值只能在内部出现。若直接在半径 \(R\) 的开盘上写“取最大点”，没有保证比值达到上确界，证明便缺少关键一步。最后让 \(r\uparrow R\)，才得到所需的边界极限比较。

\ProseParagraphBreak
曲率界也需是统一界 \(K\le-\kappa<0\)。若只知道每一点 \(K<0\)，不同半径上的最优 \(\kappa\) 可能趋于0，比较式右边随之增大，不能再让盘半径趋于无穷推出密度为0。选读小节的完整负曲率平面例子正说明这一步不能省略。Ahlfors--Schwarz 是定量的曲率比较，而不是仅由曲率符号决定全部全局几何。

\subsection{Schwarz–Pick 推广}
\label{b3:subsec:140202}

\begin{corollary}\label{b3:cor:14-schwarz-negative-target}
设 \(f:\mathbb D\rightarrow Y\) 全纯，目标有正 \(C^2\) 共形密度 \(\rho\) 且曲率不超过 \(-\kappa\)。则
\[
 \rho(f(z))|f'(z)|\le\frac{2}{\sqrt\kappa(1-|z|^2)}.
\]
\end{corollary}
\begin{proof}
若 \(f\) 常值，左端为0。否则其导数零点离散，拉回密度 \(h\) 连续，零点处为0；非零处满足同一曲率不等式。比较 \(h/\eta_r\) 的正最大值若大于1，既不在边界也不在零点，故完全适用前一定理的 Laplace 论证，得到相同上界。
\end{proof}
取目标圆盘 Poincaré 密度、\(\kappa=1\)，恢复 Schwarz--Pick 导数式。曲率方法并不要求目标的显式自同构；只要有合适密度，便可控制所有源映射。

\subsection{完备度量形式}
\label{b3:subsec:140203}

\begin{theorem}\label{b3:thm:14-hyperbolic-density-maximal}
若 \(X\) 为双曲曲面，\(\rho\) 是其任一正 \(C^2\) 共形密度且曲率不超过 \(-\kappa\)，则 \(\rho\le\lambda_X/\sqrt\kappa\)。因此对 \(f:X\rightarrow Y\) 及前一推论的目标密度，有 \(\rho(f)|f'|\le\lambda_X/\sqrt\kappa\)。
\end{theorem}
\begin{proof}
将密度沿圆盘万有覆盖拉回，对圆盘使用 Ahlfors--Schwarz，再由 Poincaré 密度下降公式消去覆盖导数，给第一式。对 \(f\) 的拉回密度作同样比较给第二式，临界点照前一推论处理。
\end{proof}
这是以完备 Poincaré 度量为源的比较形式。完备性保障后续距离几何，但本证明不需要目标密度完备。若源只是任意完备度量，尚须额外的曲率下界与全局最大值工具；不能把本定理替换成“完备便足够”的泛化。

% !TeX root = ../../Book3.tex
\section{正规族与省略值}
\label{b3:sec:1403}

省略值将像限制在一个带穿孔的球面。正规性可以通过球面导数刻画；若正规性失败，重标度会制造一条整个平面上的非恒常极限曲线，与穿孔球面的圆盘覆盖矛盾。

\subsection{省略值}
\label{b3:subsec:140301}

亚纯函数省略一个球面值，指它从不取到该值，极点被解释为取无穷远。省略三个不同球面值的族可用同一个 Möbius 变换把这三个值化为0、1、无穷远。全纯函数本已省略无穷远，所以再省略两个不同有限值便属于此情形。

\ProseParagraphBreak
球面导数定义为 \(f^\#=|f'|/(1+|f|^2)\)，极点处用 \(1/f\) 同式定义，两个表达一致。它测量到球面长度元 \(|\mathrm{d}w|/(1+|w|^2)\) 的导数大小，不是普通导数。\(f_n=n\) 的普通值无界，但球面导数0且球面趋于无穷远。

\subsection{正规性}
\label{b3:subsec:140302}

\term[marty-panbie]{Marty 判别}[Marty's criterion]以球面导数表达正规性。

\begin{theorem}[Marty 判别]\label{b3:thm:14-marty}
亚纯函数族球面正规，当且仅当其球面导数在每个紧集有共同上界。
\end{theorem}
\begin{proof}
若导数局部共同有界，球面两点的弦距离不超过连接像路径的球面长度，故在小凸盘上得到共同 Lipschitz 界。球面紧，对可数稠密点对角抽取，再用等度连续性与有限网格得到紧集一致子列。连续极限在每个点可选一个避开无穷远的球面坐标；充分晚成员也在此坐标有界全纯，Weierstrass 定理使极限全纯。所以极限为亚纯或球面常值，得到正规性。

反向若球面导数在紧集无界，取 \(f_n\)、\(z_n\rightarrow a\) 使 \(f_n^\#(z_n)\rightarrow\infty\)。正规性给球面一致子列。在 \(a\) 的极限值附近选一个球面旋转坐标，把该值送到有限点；在小盘中成员普通一致收敛且共同有界，Cauchy 导数估计给共同导数界。球面旋转保持球面导数，产生矛盾。极限恒为无穷远时同样选倒数坐标。
\end{proof}

\ProseParagraphBreak
Marty 判别把欧氏导数改成球面导数，能在极点处继续使用。若 \(g=1/f\)，则 \(|g'|/(1+|g|^2)=|f'|/(1+|f|^2)\)，两张球面坐标给相同数值。以 \(f_n(z)=nz\) 为例，0处球面导数为 \(n\)，族在那里不正规；但在任一避开0的紧集上，\(f_n\) 球面趋于无穷，正规性完全可以恢复。正规是局部性质，不能因为某一点失败就忽略其他子域的正常极限。

\ProseParagraphBreak
另一方面，\(f_n(z)=n\) 是处处球面导数为0的正规族，虽没有有限的 Euclid 局部有界性。第四章有限全纯极限的定义与此处球面正规定义使用不同紧目标，必须分清。把极限 \(\infty\) 加入球面以后，常数逃逸成为紧致现象；却没有因此把任何无界族都变成正规，\(nz\) 已给反例。

\subsection{紧性}
\label{b3:subsec:140303}

\term[zalcman-chongsuofang-yinli]{Zalcman 重缩放引理}[Zalcman's rescaling lemma]把不正规性转成非常值整平面极限。

\begin{lemma}[Zalcman 重缩放引理]\label{b3:lem:14-zalcman-rescaling}
若亚纯族在 \(a\) 附近不球面正规，则存在成员 \(f_n\)、\(a_n\rightarrow a\)、\(r_n\downarrow0\)，使 \(f_n(a_n+r_n\zeta)\) 在复平面紧集上球面趋于一个非恒常亚纯函数。
\end{lemma}
\begin{proof}
由 Marty，选择缩小圆盘半径 \(s_n\downarrow0\) 及成员，使 \(|z-a|<s_n/2\) 内某点的球面导数超过 \(n^2/s_n\)。在闭盘上最大化 \((s_n-|z-a|)f_n^\#(z)\)，取最大点 \(a_n\)，记 \(d_n=s_n-|a_n-a|\)、\(r_n=1/f_n^\#(a_n)\)。则 \(d_n/r_n\rightarrow\infty\)，且 \(r_n\rightarrow0\)。当 \(|\zeta|\le d_n/(2r_n)\)，最大性给重标度函数球面导数不超过2，而在0恰为1。对越来越大的圆盘用 Marty 与对角抽取，得到全平面亚纯极限。局部有限球面坐标中的 Cauchy 导数收敛使极限在0球面导数为1，故不恒常。
\end{proof}
重标度既选中心又选尺度；只令 \(f_n(r_nz)\) 而不寻找导数最大处，极限可能常值，无法反映正规性失败。

\ProseParagraphBreak
重缩放的目的并非单纯放大图形，而是同时固定一个非零导数和越来越大的定义域。\(f_n(z)=nz\) 在0处取尺度 \(1/n\)，得到 \(g_n(\zeta)=\zeta\)，定义域半径扩大 \(n\) 倍，极限为非常值整函数。\(f_n(z)=e^{nz}\) 在0处球面导数为 \(n/2\)，取尺度 \(2/n\) 得 \(g_n=e^{2\zeta}\)，极限仍非常值，省略0与无穷远而不省略第三个球面值。这个例子也显示三省略值定理中的数目尖锐。

\ProseParagraphBreak
如果仅固定原点而族的坏行为发生在移动点，那里导数可能不反映最大尺度。Zalcman 选择点距边界乘球面导数的最大值，保证新中心不会太靠近旧边界，并使缩放后的导数在每个固定圆盘受统一界控制。三个任务分别是捕获坏行为、扩大定义域和取得紧极限，缺少任何一个都不能完成矛盾论证。

\subsection{广义 Montel 定理}
\label{b3:subsec:140304}

\begin{theorem}[省略值 Montel 定理]\label{b3:thm:14-montel-omitted}
省略同三个不同球面值的亚纯函数族球面正规。省略同两个不同有限值的全纯函数族也球面正规，其极限为全纯函数或恒为无穷远。
\end{theorem}
\begin{proof}
若不正规，重标度给非恒常全平面亚纯极限 \(g\)。对每个省略值选局部球面坐标，Hurwitz 无零定理说明极限要么仍省略该值，要么恒取该值；非恒常排除后者。因此 \(g:\mathbb C\rightarrow\widehat{\mathbb C\kern0pt}\setminus\{a,b,c\}\)。三穿孔球面由第13章覆盖输入有圆盘万有覆盖，单连通平面使 \(g\) 提升为全纯 \(G:\mathbb C\rightarrow\mathbb D\)。Liouville 使 \(G\) 恒常，进而 \(g\) 恒常，矛盾。全纯列的球面极限若非恒为无穷远，不会产生内点极点，已由第4章 Hurwitz 倒数证明。
\end{proof}
这里补全第4章预告的省略值版本，没有在统一化或 Riemann 极值证明中回用它，因而依赖路线无循环。单个有限省略值不足，\(e^{nz}\) 在含0的盘给反例。

\ProseParagraphBreak
省略值可以随函数变化，但须避免相互碰撞。若每个 \(f_n\) 省略三个球面点 \(a_n,b_n,c_n\)，且它们两两球面距离有统一正下界，则先取子列使三点收敛到三个不同点。将它们送到 \(0,1,\infty\) 的 Möbius 变换及其逆随参数连续，在球面上一致收敛到非退化变换；正规三省略值定理应用于变换后的族，再复合逆变换即可。此论证也给族的任意序列子列判别。

\ProseParagraphBreak
若允许碰撞，结论失败。例如 \(f_n(z)=nz\) 在单位盘省略 \(2n,3n,4n\)，三个点球面上都趋向无穷，族仍在0不正规。“每个函数都有三个省略值”若不记录这三个值共同固定或有统一分离，不能当作 Montel 定理的假设。

% !TeX root = ../../Book3.tex
\section{Picard 小定理}
\label{b3:sec:1404}

Picard 小定理把省略值从族的紧性转成单个整函数的刚性。圆盘覆盖将目标限制变成一个有界整函数，Liouville 遂可直接使用。

\subsection{覆盖空间证明}
\label{b3:subsec:140401}

\term[picard-xiao-dingli]{Picard 小定理}[little Picard theorem]限制整函数的省略值。

\begin{theorem}[Picard 小定理]\label{b3:thm:14-picard-small}
非恒常整函数至多省略一个有限复值；非恒常全平面亚纯函数至多省略两个球面值。
\end{theorem}
\begin{proof}
若亚纯函数省略三个不同球面值，它是从单连通平面到三穿孔球面的全纯映射。路径提升给到圆盘万有覆盖的全纯提升，Liouville 使提升常值，所以原函数常值。整函数本来不取球面值无穷远，再省略两个有限值就省略三个球面值，因此也不可能。
\end{proof}
这项证明使用已声明的深层覆盖模型输入，提升与 Liouville 部分完整。它不将“三穿孔球面有圆盘覆盖”从 Picard 定理反推，避免循环。

\subsection{省略值}
\label{b3:subsec:140402}

\(e^z\) 省略0，说明整函数的一值例外确会发生。作为亚纯函数它同时省略0与无穷远，达到球面版本的两值上界。多项式则不省略任何有限值，代数基本定理已经给出其满射性。

\ProseParagraphBreak
“至多一个”不表示必有一个。\(\sin z\) 对每个 \(a\in\mathbb C\) 都取值：令 \(w=e^{\mathrm{i}z}\)，方程 \((w-w^{-1})/(2\mathrm{i})=a\) 化为常数项为 \(-1\) 的二次方程，根非零，再取复对数求 \(z\)。小定理只控制遗漏，不直接计数原像。

\subsection{应用}
\label{b3:subsec:140403}

若非常值整函数 \(f,g\) 满足 \(e^f+e^g=1\)，则 \(e^f\) 同时不取0与1，Picard 小定理迫使它常值，求导再给 \(f\) 常值，与假设矛盾。因此不存在这样的非常值整函数对。相同方法可检验一些指数型函数方程：先把式子变成两个省略值，再用刚性，而非展开无穷系数逐项比较。

\ProseParagraphBreak
区域换成圆盘后这种结论就不成立。常值附近足够小的全纯扰动可同时避开0与1。整平面的尺度无限是 Liouville 提升证明中的决定条件，不能把小定理搬到任意区域。

\ProseParagraphBreak
对于超越整函数，无穷远是本性奇点。理由是把 \(F(z)\) 改写成 \(F(1/t)\)：若在0可去，则 \(F\) 在大圆外有界，再与圆内紧性合并，由 Liouville 为常数；若为极点，则 \(F\) 有多项式增长，Cauchy 导数估计迫使充分高阶 Taylor 系数为0，故为多项式。超越整函数既非前者也非后者，所以大 Picard 给它在每个大圆外取遍所有复值至多除去一个，而且其余值取无限次。

\ProseParagraphBreak
小 Picard 只限制全局省略值，大 Picard 还控制重复次数。\(e^z\) 省略0，其他每个非零值在任意大圆外都有无穷多个对数原像。\(\sin z\) 不省略任何复值，利用 \(e^{\mathrm{i}z}\) 的二次方程及周期也可显式看到无限原像。这些例子与定理一致，但不能由几个常见超越函数的行为代替一般本性奇点的证明。

% !TeX root = ../../Book3.tex
\section{Picard 大定理}
\label{b3:sec:1405}

本性奇点附近的稠密像由 Casorati--Weierstrass 给出，Picard 大定理进一步保证除至多一值外，每个值在任意穿孔邻域取到无穷次。正规族把局部遗漏转成整平面极限，是增强结论的桥梁。

\subsection{本性奇点}
\label{b3:subsec:140501}

设 \(f\) 在 \(0<|z-a|<R\) 全纯且 \(a\) 为本性奇点。这里本性是 Laurent 主部无限，不是简单的无界性；极点同样无界，却在小邻域遗漏一个大圆盘中的全部值。

\ProseParagraphBreak
若某个值在一个穿孔邻域只有有限个原像，缩小邻域后便可完全遗漏该值。因此证明大定理只需证明本性奇点附近不能同时遗漏两个不同有限值，再把这个结论用于有限原像的两个候选例外。

\subsection{Picard 大定理}
\label{b3:subsec:140502}

\term[picard-da-dingli]{Picard 大定理}[great Picard theorem]限制本性奇点附近的例外值。

\begin{theorem}[Picard 大定理]\label{b3:thm:14-picard-great}
在全纯函数的本性孤立奇点附近，除至多一个有限值外，每个复值在任意穿孔邻域中取得无穷次。
\end{theorem}
\begin{proof}
平移奇点到0。假设某穿孔盘上遗漏两个值，仿射归一化为0与1。选 \(b\notin\{0,1\}\)，Casorati--Weierstrass 给 \(z_n\rightarrow0\) 且 \(f(z_n)\rightarrow b\)。重缩放 \(g_n(\zeta)=f(z_n\zeta)\) 在 \(0<|\zeta|<R/|z_n|\) 定义，遗漏0、1、无穷远。对 \(\mathbb C^*\) 的紧穷竭用省略值 Montel 和对角抽取，得到球面极限 \(g\)，且 \(g(1)=b\)。

若 \(g\) 非恒常，它仍遗漏三值；复合指数 \(g(e^w)\) 是全平面非恒常亚纯函数且遗漏三值，与 Picard 小定理矛盾。故 \(g\equiv b\)。由于这个常数有限，收敛在单位圆上是普通一致，充分晚有 \(|f(z_n\zeta)|\le M\)、\(|\zeta|=1\)。抽取 \(|z_n|\) 严格递减子列，这给趋于0的一列完整圆周上的同一界。相邻圆周之间 \(f\) 全纯且连续到边界，最大模原理使整个闭圆环有同一界。这些圆环覆盖一个穿孔邻域，所以 \(f\) 有界，0可去，矛盾。

因此两个遗漏值不可能。若两个不同值各只有有限次原像，缩小穿孔盘同时避开这些有限原像，便得到两个遗漏值，同样矛盾，故至多一值例外。
\end{proof}
对亚纯本性奇点，至多两个球面值例外；若有三个，先把其中一个送到无穷远，使小穿孔盘上函数全纯，再用刚证的版本。

\subsection{与 Casorati–Weierstrass 定理的比较}
\label{b3:subsec:140503}

\(e^{1/z}\) 在0本性且遗漏0；对每个 \(b\ne0\)，解为 \(z=(\log b+2\pi\mathrm{i}n)^{-1}\)，除去可能分母为0的一项，得到无穷原像趋于0。这个例子使全部量词都可直接核对：每个非零值、每个穿孔邻域、无穷多个原像。

\ProseParagraphBreak
Casorati--Weierstrass 只保证像稠密，从接近某值不能推出真的取到该值。大定理的证明通过正规族、覆盖与最大模原理填补这个差距，不把“稠密”误写成“满射”。大定理也不声称例外值必存在，某些本性函数没有有限例外。

\ProseParagraphBreak
大 Picard 中的“无限次”可从缩小邻域的形式严格推出。若某值 \(a\) 在一个穿孔邻域只有有限个原像，就把圆盘缩到排除全部这些原像；它便在更小穿孔邻域省略 \(a\)。若有两个不同值各只取有限次，缩到同时排除两组原像便得到两个省略值，违反全纯大 Picard。亚纯版本同理，但容许两个球面例外值而非一个有限例外值。

\ProseParagraphBreak
Casorati--Weierstrass 的稠密像不保证每个指定值真的出现。一个开稠密集仍可以缺少无限多个点；本性奇点的稠密性只说目标任何小邻域都有像。Picard 则把未取值的集合压缩到至多一个或两个点。证明中利用稠密性选择一个通常值的逼近序列，再利用正规族把局部逼近转成一整圈的统一控制，最后用最大模原理推出可去性。稠密性、紧致性和环带估计各自承担不同步骤，缺少其中之一就不能从稠密像升级到省略值结论。

% !TeX root = ../../Book3.tex
\OptionalSection{曲率方法}{b3:sec:1406}

本选读节说明负曲率与重标度如何组织更广的全纯刚性。首先保留统一曲率常数的准确范围，再以紧目标讨论整个曲线与正规性的联系。

\subsection{\texorpdfstring{\(\mathbb C\)}{C} 上不存在负曲率度量}
\label{b3:subsec:140601}

\begin{corollary}\label{b3:cor:14-no-uniform-negative-plane}
全平面不存在正 \(C^2\) 共形密度满足 \(K\le-\kappa<0\)。
\end{corollary}
\begin{proof}
对任意中心 \(a\) 的半径 \(R\) 圆盘使用 Ahlfors--Schwarz，在中心得到 \(\rho(a)\le2/(\sqrt\kappa R)\)。令 \(R\rightarrow\infty\)，得到 \(\rho(a)=0\)，违反正性。
\end{proof}
这里只排除统一负上界，不排除处处负曲率。\(\rho=e^{|z|^2}\) 给 \(K=-4e^{-2|z|^2}<0\)，且度量完备，提供准确反例。标题中的“负曲率”应按前述量化陈述理解，不能删除 \(\kappa\) 而使定理变成错误命题。

\subsection{Liouville 型结论}
\label{b3:subsec:140602}

若目标有统一负曲率密度，任意整平面全纯映射到它都常值。把映射限制在以任意源点为中心的大圆盘，导数比较给 \(\rho(f(a))|f'(a)|\le2/(\sqrt\kappa R)\)；令 \(R\rightarrow\infty\)，目标密度正使 \(f'(a)=0\)。这是一种几何 Liouville 结论，目标未必在某个坐标中有界。

\ProseParagraphBreak
对于双曲曲面，Poincaré 密度曲率 \(-1\)，所以整曲线均常值；也可直接提升到圆盘使用 Liouville。两个证明分别突出曲率与覆盖，所需前提等价于已建立的双曲模型。

\subsection{Brody 型思想}
\label{b3:subsec:140603}

重标度还适用于紧 Riemann 曲面目标的任一光滑 Hermitian 度量。若圆盘映射列导数在内点附近无共同界，按本章最大化“到边界距离乘导数”的方法选中心与尺度，得到越来越大圆盘上导数有共同界、中心导数为1的映射列。目标紧给取值紧性，等度连续性给全平面极限；在局部复坐标中用 Cauchy 公式使极限全纯且中心导数非零。因此出现非恒常整曲线。

\ProseParagraphBreak
反向，若有非恒常整曲线 \(g:\mathbb C\rightarrow X\)，在某个非零导数点平移后考察 \(g(nz)\)，其中心导数无界，故圆盘映射族不正规。这是曲面情形的 Brody 重标度思想。紧性不可省略，否则像可能逃向目标外部而没有全平面极限；球面 Zalcman 证明可工作，正因为球面本身紧。

\subsection{双曲性理论的进一步讨论}
\label{b3:subsec:140604}

紧球面有非恒常整曲线，例如 \(z\mapsto z\) 视为球面值映射；复环面也有投影 \(\mathbb C\rightarrow\mathbb C/\Lambda\)。紧的圆盘型曲面则没有这样的曲线。因此全纯曲线刚性识别了紧曲面三种覆盖模型中的双曲情形。

\ProseParagraphBreak
更高维中，负曲率、Kobayashi 距离和整个曲线之间的关系需另立定义与定理，不能只把一维 Gaussian 曲率公式改几个指标后宣称全部成立。本卷以后讨论双曲性时保留一维范围；本节只提供从导数比较、紧性到全局刚性的证明思路。

\ProseParagraphBreak
本章有两种从局部估计取得全局结论的方法。曲率比较在越来越大的圆盘上控制拉回密度，极限使整映射的导数消失；重缩放则从坏的局部导数中提取一个非平凡整映射，再以省略值或目标双曲性排除它。两者所用的几何尺度不同，但都必须核对新定义域是否真的趋于整个平面，且归一化是否保留非平凡性。

\ProseParagraphBreak
在球面和环面中，非恒定整曲线确实存在：前者可取 \(z\mapsto z\) 作为球面值映射，后者取商投影 \(\mathbb C\to\mathbb C/\Lambda\)。圆盘覆盖的曲面则没有，源平面单连通使它提升到圆盘，再由 Liouville 为常数。这个区别是在本书已经标明的统一化输入上得到的完整推论；更一般复流形的 Brody 定理需要额外紧致性和几何论证，不在这里扩大声称。


\begin{chaptersummary}
统一负曲率给全纯导数上界，球面导数给正规族判别，重标度把不正规变成整个曲线。三穿孔球面的圆盘覆盖排除该曲线，得到省略值 Montel；同一覆盖输入给小 Picard，再由缩放和圆环最大模推出大 Picard。
\end{chaptersummary}

\BookExercises

% 本章题组自拟；以正文已建立的工具设计计算、反例、构造与综合训练。
\begin{exercise}\label{b3:ex:14-1}
求 \(\rho=e^{|z|^2}\) 的曲率，说明它与全平面不存在结论不矛盾。
\end{exercise}
\begin{solution}
\(\Delta\log\rho=4\)，故 \(K=-4e^{-2|z|^2}<0\)。它趋于0，没有 \(K\le-\kappa\) 的统一常数，所以不满足不存在定理的假设。
\end{solution}

\begin{exercise}\label{b3:ex:14-2}
若密度曲率为 \(-1\)，整体乘3后曲率为多少？
\end{exercise}
\begin{solution}
对数仅加常数，Laplace 不变，密度平方乘9，所以曲率为 \(-1/9\)。长度归一化与曲率常数必须一起改变。
\end{solution}

\begin{exercise}\label{b3:ex:14-3}
解释拉回度量在 \(f(z)=z^2\) 的0点为何不是正度量，比较证明为何仍可使用。
\end{exercise}
\begin{solution}
拉回密度含 \(|2z|\)，0处为0，曲率式不能直接除以其平方。若与正比较密度的比值有大于1的最大值，它必避开零点，于是在该最大点普通曲率论证有效。
\end{solution}

\begin{exercise}\label{b3:ex:14-4}
证明族 \(e^{nz}\) 在含0的盘不球面正规，指出它省略的值数。
\end{exercise}
\begin{solution}
球面导数在0为 \(n/2\)，Marty 判别给不正规。它省略0与无穷远，共两个球面值，尚未达到省略三个值的定理条件。
\end{solution}

\begin{exercise}\label{b3:ex:14-5}
证明若整函数 \(f\) 同时使 \(f\) 与 \(f-1\) 无零，则 \(f\) 常值。
\end{exercise}
\begin{solution}
它遗漏两个有限值0、1；作为球面值函数再遗漏无穷远，Picard 小定理给常值。条件对圆盘函数则不够。
\end{solution}

\begin{exercise}\label{b3:ex:14-6}
对 \(e^{1/z}\)，求非零值 \(b\) 的原像并说明为何任意穿孔盘有无穷多个。
\end{exercise}
\begin{solution}
原像 \(z_n=(\log b+2\pi\mathrm{i}n)^{-1}\)，去掉至多一个零分母项。其模趋于0，点互不相同，因此任意穿孔邻域包含无穷多个。0为唯一有限例外。
\end{solution}

\begin{exercise}\label{b3:ex:14-7}
解释大 Picard 证明中“完整圆周上有共同界”为什么足以得到整个穿孔盘有界。
\end{exercise}
\begin{solution}
把趋于0的圆周半径抽成严格递减列。相邻圆周形成闭圆环，函数在邻域全纯，由边界最大模原理同界控制内部。这些圆环覆盖全部充分小非零点，所以得全穿孔界。只沿一条径向序列有界则不足。
\end{solution}

\begin{exercise}\label{b3:ex:14-8}
证明从全平面到任一双曲曲面的全纯映射恒常，并说明复环面为何不适用。
\end{exercise}
\begin{solution}
源平面单连通，映射提升到目标圆盘万有覆盖，提升有界整，由 Liouville 常值。复环面覆盖是平面，存在非恒常整个投影，因此不是双曲目标。
\end{solution}
